\subsubsection{Repeatability and scope choices}
Fifty five of sixty repeated judgments matched. The five changes are on A1, A3 and A5, and three of them concern A3, where the responses weigh partition statements and edition metadata differently between repetitions. For HumanEval B4, both inputs and both repetitions use 5 because the paper compares expected student performance with model results in its discussion. The Codex batch had used not applicable for Input 1. For DesignQA B4, all four Claude responses use not applicable because no human performance is measured. These scope choices are stable within this batch, but the rubric does not say whether an interpretive comparison in a discussion section brings B4 into scope, so stability is not evidence that the choice is right.

One evidence entry in a DesignQA Input 2 response cites the file listing itself as the source of an absence claim about a license file. The citation checker counts it as a source that does not exist, because the listing is a tool result rather than a file. The response states the inference plainly, and the snapshot has no license file. This shows that the mechanical source check and the support for a judgment can diverge in both directions.

The observations above cover all sixty judgments in the six preselected responses through the accompanying audit. The analyst is the same model as the reviewer, and no independent human analysis was performed. No semantic accuracy rate is inferred from quotation matching, from agreement between the batches or from the selected examples.
